\subsection{A local insertion into the prescribed simple-torus family}

\begin{proposition}[B.4.15; Proved]\label{push:B415}
Let $q_0$ be a regular point of the elliptic base, and let
$\Delta_z$, $z=q-q_0$, be a simply connected disc avoiding
$0,1,\infty$. After $w=r^6$, take the standard twelve-cycle
model
\[
 \mathcal X\longrightarrow\Delta_z\times\Delta_r
\]
of the prescribed local two-torus family. There is a proper
replacement $\mathcal X'\to\Delta_z\times\Delta_r$, with
a specified isomorphism to $\mathcal X$ off $(z,r)=(0,0)$,
whose total fourfold is smooth and whose divisor $r=0$ is
reduced normal crossings. The composite map to $\Delta_r$
is flat.

The replacement has twelve new compact vertical normalization
components. Six are $\PP^1$-bundles over a scalar Hopf
surface; the other six are the blow-ups of those bundles
along two disjoint elliptic curves. Every new component
has fundamental group $\Z$, and therefore fails the
$\partial\bar\partial$ lemma. The punctured family is the
actual prescribed family, so its generic simplicity is
preserved.

This is a statement about a proper model over a bidisc,
flat over the degeneration parameter. It does not assert
joint flatness over the bidisc, a one-parameter degeneration
of compact surfaces with Hopf components, or a global
$S^6$ filling. All twelve component labels, and hence the
obstruction of B.3.1, are unchanged.
\end{proposition}

\begin{proof}
\emph{The actual action.}
Use the rank-three period subsystem
\[
 \mathbf u(z)=(\tau(z),\mu(z)),\qquad
 \mathbf w=(1,0),\qquad \boldsymbol\delta=(0,1).
\]
It defines the analytic semiabelian group
\[
 G_z=\C^2/(\Z\mathbf u(z)+\Z\mathbf w+
                    \Z\boldsymbol\delta),
 \qquad
 1\longrightarrow\C^*_{\delta}\longrightarrow G_z
       \longrightarrow E_z\longrightarrow0,
 \quad E_z=\C/(\Z+\tau(z)\Z).
\]
This is the punctured degree-zero line bundle determined
by the actual $\mu(z)$. We make neither a product assumption
nor a torsion assumption on that line bundle. The formula
\[
 \chi_z(t)=
 \left[\frac{\log t}{2\pi i}\,\mathbf u(z)\right]
 \quad(t\in\C^*)
 \tag{B.4.15.1}
\]
is a single-valued holomorphic homomorphism: a logarithm
change adds an actual lattice period. It varies
holomorphically with $z$. The $r$-monodromy fixes
$\mathbf u$, and the $z$-disc has no period monodromy.

Here the action extends to the entire twelve-cycle model,
not just to its smooth fibres. Each unfolded normalization
is
\[
 R_z=G_z\mathbin{\times^{\C^*_{\delta}}}\PP^1
       =\PP_{E_z}(\OO\oplus L_z).
\]
Translation in $G_z$ preserves both conductor sections.
More explicitly, a class represented by $(a,b)$ acts
before the closing quotient by
\[
  (\xi,W)\longmapsto
            (\xi+a,e^{-2\pi i b}W).
 \tag{B.4.15.2}
\]
It commutes exactly with the closing period, which is
another translation of the same commutative group and
whose scalar factor contains $r^{12}$. The adjacent line
directions are paired dually in the local equation $xy=r$;
the two scalar factors cancel. Scalar multiplication also
preserves the one-dimensional toric subdivision producing
the twelve components. Thus (B.4.15.2) extends over every
node, preserves each branch, and changes a branch equation
$x$ by a unit times $x$. This checks the required action
using the actual semiabelian marking of B.1.

The affine family itself admits the necessary local
punctured comparison. On
$U=\Delta_z\times\Delta_r^*$ its lattice local system
$\Lambda$ has $H^2(U,\Lambda)=0$, since $U$ retracts onto
a circle, also for local coefficients. Its relative Lie
algebra $\mathcal V$ is a coherent vector bundle. The
exponential sequence
\[
 0\longrightarrow\Lambda\longrightarrow\mathcal V
       \longrightarrow\mathcal T\longrightarrow0
\]
and Cartan's theorem B on the Stein base give
$H^1(U,\mathcal T)=0$ \citep{Cartan1952B_AB}.
Consequently the actual affine torsor has a holomorphic
section on $U$, which supplies the stated punctured
comparison with its Jacobian. No extension of this
section across $r=0$ is asserted.

\emph{The base blow-up and small resolutions.}
Blow up the base origin:
\[
 \beta:\widetilde B=\operatorname{Bl}_{(0,0)}
            (\Delta_z\times\Delta_r)
       \longrightarrow\Delta_z\times\Delta_r,
 \qquad F=\beta^{-1}(0,0)\cong\PP^1.
\]
In the chart $r=zs$, the pulled-back node is
$xy=zs$, with one further smooth elliptic coordinate.
Label the original components by $i\in\Z/12$.
At each node choose $x=0$ as the even-indexed branch,
and blow up $(x,z)$. This is globally the sum of the
pulled-back invertible ideal of the chosen central
component and the ideal $\OO_{\widetilde B}(-F)$.
The ideals are $G$-invariant. Their non-Cartier loci are
the twelve disjoint elliptic node loci, so the choices
at the six even components are compatible.

The two strict-transform charts are
\[
 \begin{aligned}
 z&=xQ,&y&=Qs,&r&=xQs,\\
 x&=zX,&s&=Xy,&r&=zXy.
 \end{aligned}
 \tag{B.4.15.3}
\]
They are smooth, and all three factors of $r$ occur once.
In the other base chart $z=r\zeta$, the node is $xy=r$
and the same ideal is $(x,r)=(x)$, so no change occurs.
Away from the indicated node loci it is also invertible.
The resulting total space is smooth and its $r$-central
divisor is reduced normal crossings.

For the component identifications, before resolution put
\[
 \begin{array}{ll}
 V_x=\{x=z=0\},&V_y=\{y=z=0\},\\
 H_x=\{x=s=0\},&H_y=\{y=s=0\}.
 \end{array}
\]
The $V$ branches are vertical over $F$, and the $H$
branches lie over the strict transform of $r=0$.
On $V_x$, the two maps in (B.4.15.3) are
\[
 (Q,s)\longmapsto(y,s)=(Qs,s),\qquad
 (X,y)\longmapsto(y,s)=(y,Xy).
\]
These are exactly the blow-up charts of $(y,s)$.
Thus $V_x$ is blown up along its elliptic curve over
$y=s=0$; $V_y$ is unchanged. Similarly $H_y$ is blown
up along the corresponding elliptic curve and $H_x$
is unchanged. Each even vertical component is therefore
blown up at its two disjoint conductor curves, whereas
each odd vertical component is unchanged. The horizontal
parities are reversed.

\emph{The Cartier logarithm and properness.}
Set $M=\OO_{\widetilde B}(F)$. On a trivializing cover
write its transition units as $g_{ij}$. Glue the resolved
pullbacks by $\chi_z(g_{ij})$. These are genuine
automorphisms satisfying the cocycle identity.
This is the Cartier-logarithm construction of
\citet[Definition~1.1]{Fujimoto1990AB}; the explicit
semiabelian action above verifies its applicability
here, without assuming that the smooth two-tori are
principal elliptic bundles. The invariant ideals show
that the action lifts to the resolutions; equivalently
one may twist first and then resolve the descended ideals.

Over every base chart the glued space is a copy of the
proper resolved pullback. It is Hausdorff: two points
over the same base point are separated in such a chart,
and points over distinct base points are separated in
the base. Properness is local on the target, and
composition with the proper $\beta$ proves properness
over the original bidisc. The canonical divisor section
of $M$ is nonzero off $F$ and trivializes the twist there.
The base blow-up and the resolutions are also
isomorphisms off the origin. This gives the specified
isomorphism with the original family there, in particular
over every $r\ne0$.

The construction can be inserted into any larger model
which is standard near $q_0$: the preceding
trivialization supplies the overlap map outside the
origin. The cocharacter $\mathbf u$ is needed only on
this disc and is not extended around the special base
points. Since the local rings of the total space are
regular and $r$ is a nonzero divisor, they are
torsion-free, hence flat, over the one-dimensional
regular parameter ring. This proves flatness over $r$.
The new vertical threefolds rule out joint flatness over
the bidisc. Properness over the $r$-disc alone follows
only when this local replacement is inserted into a
model with compact elliptic base.

\emph{The compact vertical components.}
Before resolving, each vertical normalization is the
twist of $R_{q_0}\times F$ by
$M|_F=\OO_{\PP^1}(-1)$, namely
\[
 V_i^0=\OO_{\PP^1}(-1)^\times
                    \mathbin{\times^{\C^*}}R_{q_0}.
\]
Its elliptic quotient is a principal $E_{q_0}$-bundle
$H\to\PP^1$ with primitive Chern cycle
$\pm\mathbf u$. It is explicitly a scalar Hopf surface.
Indeed put
\[
 \tau_0=\tau(q_0),\qquad
 \eta=-1/\tau_0,\qquad a'=e^{2\pi i\eta},
 \qquad |a'|<1.
\]
After dividing the elliptic coordinate by $\tau_0$,
the projection of (B.4.15.1) is
$[\log t/(2\pi i)]$ in $\C/(\Z+\eta\Z)$.
Its kernel is $(a')^\Z$. Consequently
\[
 H=\OO_{\PP^1}(-1)^\times/(a')^\Z
   =(\C^2\setminus\{0\})/\langle a'\id\rangle.
 \tag{B.4.15.4}
\]
The principal $G$-bundle
$\OO(-1)^\times\mathbin{\times^{\C^*}}G_{q_0}$
maps to $H$ with fibre $\C^*_{\delta}$.
Compactifying this principal $\C^*$-bundle gives an
actual projective bundle
\[
 V_i^0=\PP_H(\OO\oplus M_H)
\]
with its two global conductor sections. On each
elliptic fibre of $H\to\PP^1$, the line $M_H$ retains
the original degree-zero line $L_{q_0}$.

The two conductor sections over the distinguished point
$F\cap\operatorname{strict}(r=0)$ are disjoint elliptic
curves. The component calculation in (B.4.15.3) gives
exactly the six unchanged bundles and six two-curve
blow-ups in the statement. A $\PP^1$-bundle has the
fundamental group of its base, and blowing up a smooth
complex codimension-two centre preserves fundamental
group. Hence every new normalization has
$\pi_1=\pi_1(H)=\Z$. A compact
$\partial\bar\partial$ manifold has even first Betti
number, so these normalizations fail that lemma.

Finally, the connected translation action fixes every
component label. The rotations $+4,-3,0$ of the
incompatible global assembly are unchanged. This proves
the precise local assertion and its stated limitation.
\end{proof}

\begin{remark}
For a generic arc $z=kr$ through the blown-up base
origin, the strict transform meets $F$ at an ordinary
point. Its central surface remains a cycle of ruled
surfaces over the elliptic curve. The non-K\"ahler
geometry in (B.4.15.4) is carried by the variation
around the whole exceptional $\PP^1$. Therefore a
theorem only about one-parameter limits of simple
two-tori would not, by itself, exclude these vertical
threefold components.
\end{remark}
